\providecommand{\todo}[1]{\textbf{[TODO: #1]}}
% ============================
% ภาคผนวก
% ============================
\setcounter{chapter}{-1}

\appendiceschapter{การตั้งค่าระบบ}
\section*{คำสั่งประมวลผลเสียง}

ตัวกรองเสียงที่ใช้งานจริงสำหรับการประมวลผลสัญญาณเสียงเบื้องต้นผ่าน FFmpeg แสดงดังนี้ โดยการแปลงความถี่เป็น 16 kHz แบบ Mono ดำเนินการผ่านโมดูลโหลดสัญญาณเสียงของไลบรารีถอดความ:

\begin{verbatim}
ffmpeg -i input.mp3 -af "afftdn,
  silenceremove=start_periods=1:start_duration=0.1:
  start_threshold=-40dB:stop_periods=-1:
  stop_duration=0.6:stop_threshold=-40dB" output.wav
\end{verbatim}

\begin{table}[H]
\centering
\caption{พารามิเตอร์ของตัวกรองเสียง}
\label{tab:appfilter}
\small
\begin{tabular}{|l|l|p{6.0cm}|}
\hline
\textbf{ตัวกรอง} & \textbf{พารามิเตอร์} & \textbf{ค่า} \\
\hline
afftdn & nr / nf / nt & 12 dB / $-50$ dB / white (ค่าเริ่มต้นของ FFmpeg) \\
\hline
silenceremove & ต้นไฟล์ & ตัดเสียงเงียบที่ยาวเกิน 0.1 วินาที ต่ำกว่า $-40$ dB \\
\hline
silenceremove & ตลอดไฟล์ & ตัดช่วงเงียบที่ยาวเกิน 0.6 วินาที ต่ำกว่า $-40$ dB \\
\hline
\end{tabular}
\end{table}

\section*{ค่าตั้งค่าการถอดความ}

\begin{table}[H]
\centering
\caption{ค่าตั้งค่าของเอนจินถอดความ}
\label{tab:appconfig}
\small
\begin{tabular}{|l|l|}
\hline
\textbf{พารามิเตอร์} & \textbf{ค่า} \\
\hline
แบบจำลอง & Faster-Whisper Small \\
\hline
อุปกรณ์ & CPU, \texttt{cpu\_threads=8} \\
\hline
Compute type & INT8 \\
\hline
Beam size & 1 \\
\hline
\end{tabular}
\end{table}

\section*{ข้อความนำ (PATHOLOGY\_PROMPT)}

\begin{quote}
\small
Surgical pathology gross examination report. Surgical number S-26-1001. Received in formalin is a right modified radical mastectomy specimen. Left simple mastectomy. Skin ellipse. The nipple is everted, inverted, unremarkable, shows ulceration. Infiltrative firm yellow-white mass located at upper outer quadrant, upper inner quadrant, lower outer quadrant, lower inner quadrant, central, subareolar. Well-defined firm white mass with slit-like appearance. Poorly circumscribed yellow-white lesion. Previous surgical cavity with adjacent fibrous tissue. Residual mass. Deep surgical resection margin, superior margin, inferior margin, medial margin, lateral margin, skin margin. Centimeters, cm, millimeters, mm. Distance from closest margin. Grossly free from tumor. Representative sections submitted in paraffin blocks. Axillary contents, level 1, level 2, lymph nodes.
\end{quote}

\noindent ข้อความนี้มี 111 คำ และเมื่อนับด้วยตัวตัดคำของ Whisper แบบหลายภาษาได้ 198 โทเคน ซึ่งไม่เกินขีดจำกัดของข้อความนำที่ 223 โทเคน

\appendiceschapter{ฟิลด์ข้อมูลที่สกัดและผลการประเมินโดยละเอียด}
\section*{รายการฟิลด์ 15 ฟิลด์}

\begin{table}[H]
\centering
\caption{ฟิลด์ที่ตัวสกัดข้อมูลรองรับ}
\label{tab:appfields}
\footnotesize\setlength{\tabcolsep}{3pt}
\begin{tabular}{|c|p{5.0cm}|p{6.6cm}|}
\hline
\textbf{ที่} & \textbf{รหัสฟิลด์} & \textbf{ความหมาย} \\
\hline
1 & \texttt{s0\_surgical\_no} & เลขที่สิ่งส่งตรวจทางศัลยกรรม \\
\hline
2 & \texttt{s1\_side} & ด้านของเต้านม (Right / Left) \\
\hline
3 & \texttt{s2\_proc} & ประเภทหัตถการ (Mastectomy, Lumpectomy ฯลฯ) \\
\hline
4 & \texttt{s3\_dims} & ขนาด 3 มิติของชิ้นเนื้อ (ซม.) \\
\hline
5 & \texttt{s4\_skin} & มีชิ้นผิวหนังติดมาหรือไม่ \\
\hline
6 & \texttt{s5\_dims} & ขนาด 2 มิติของชิ้นผิวหนัง (ซม.) \\
\hline
7 & \texttt{s6\_nipple} & สภาพหัวนมและลานนม \\
\hline
8 & \texttt{s7\_biopsy\_scar} & มีรอยแผลเป็นจากการเจาะตรวจเดิมหรือไม่ \\
\hline
9 & \texttt{s8\_cavity} & มีโพรงจากการผ่าตัดครั้งก่อนหรือไม่ \\
\hline
10 & \texttt{s9\_residual\_mass} & พบก้อนหลงเหลือติดโพรงเดิมหรือไม่ \\
\hline
11 & \texttt{s10\_infiltrative} & ลักษณะของก้อนหลัก (พบก้อนลุกลามหรือไม่) \\
\hline
12 & \texttt{s10\_inf\_dims} & ขนาด 3 มิติของก้อนหลัก \\
\hline
13 & \texttt{s10\_5\_quadrant} & ตำแหน่งจตุภาคที่พบก้อน \\
\hline
14 & \texttt{s11\_deep\_margin} & ระยะห่างจากขอบตัดด้านลึก (Deep Margin) \\
\hline
15 & \texttt{s14\_check} & การตรวจต่อมน้ำเหลืองรักแร้ (ค่าจริงหรือเท็จ: พบหรือไม่พบต่อมน้ำเหลือง) \\
\hline
\end{tabular}
\end{table}

\noindent ฟิลด์ที่มีค่าบวกในเฉลยของชุดทดสอบ 1,000 กรณีมี 11 ฟิลด์ (ทุกฟิลด์ยกเว้น \texttt{s6\_nipple}, \texttt{s7\_biopsy\_scar}, \texttt{s8\_cavity} และ \texttt{s9\_residual\_mass}) ส่วน 4 ฟิลด์ดังกล่าวเป็นฟิลด์ควบคุมเชิงลบที่ไม่มีการระบุค่า จึงรายงานเฉพาะค่าความจำเพาะ ทั้งนี้ฟิลด์ \texttt{s14\_check} ใช้ระบุการตรวจพบต่อมน้ำเหลือง และฟิลด์ \texttt{s11\_deep\_margin} ครอบคลุมการประเมินระยะห่างของขอบตัดด้านลึกตามขอบเขตที่กำหนด

\section*{จำนวนนับต่อฟิลด์}

\begin{table}[H]
\centering
\caption{จำนวนนับ TP, FP, FN และ TN ต่อฟิลด์ (1,000 กรณี) กรณีสกัดได้ค่าผิดนับเป็นทั้ง FP และ FN}
\label{tab:confmat}
\scriptsize\setlength{\tabcolsep}{2.5pt}
\begin{tabular}{|l|r|r|r|r|r|r|r|r|}
\hline
 & \multicolumn{4}{c|}{\textbf{Baseline}} & \multicolumn{4}{c|}{\textbf{PathoWhisper}} \\
\hline
\textbf{ฟิลด์} & \textbf{TP} & \textbf{FP} & \textbf{FN} & \textbf{TN} & \textbf{TP} & \textbf{FP} & \textbf{FN} & \textbf{TN} \\
\hline
\texttt{s0\_surgical\_no} & 976 & 24 & 24 & 0 & 978 & 20 & 22 & 0 \\
\hline
\texttt{s1\_side} & 898 & 0 & 102 & 0 & 999 & 1 & 1 & 0 \\
\hline
\texttt{s2\_proc} & 900 & 0 & 0 & 100 & 900 & 0 & 0 & 100 \\
\hline
\texttt{s3\_dims} & 983 & 12 & 17 & 0 & 977 & 14 & 23 & 0 \\
\hline
\texttt{s4\_skin} & 481 & 0 & 19 & 500 & 498 & 0 & 2 & 500 \\
\hline
\texttt{s5\_dims} & 376 & 6 & 24 & 600 & 393 & 5 & 7 & 600 \\
\hline
\texttt{s6\_nipple} & 0 & 0 & 0 & 1,000 & 0 & 0 & 0 & 1,000 \\
\hline
\texttt{s7\_biopsy\_scar} & 0 & 0 & 0 & 1,000 & 0 & 0 & 0 & 1,000 \\
\hline
\texttt{s8\_cavity} & 0 & 0 & 0 & 1,000 & 0 & 0 & 0 & 1,000 \\
\hline
\texttt{s9\_residual\_mass} & 0 & 0 & 0 & 1,000 & 0 & 0 & 0 & 1,000 \\
\hline
\texttt{s10\_infiltrative} & 782 & 2 & 18 & 198 & 794 & 12 & 6 & 188 \\
\hline
\texttt{s10\_inf\_dims} & 653 & 21 & 47 & 299 & 660 & 17 & 40 & 300 \\
\hline
\texttt{s10\_5\_quadrant} & 500 & 0 & 0 & 500 & 498 & 2 & 2 & 500 \\
\hline
\texttt{s11\_deep\_margin} & 687 & 1 & 13 & 300 & 656 & 6 & 44 & 300 \\
\hline
\texttt{s14\_check} & 575 & 0 & 25 & 400 & 598 & 0 & 2 & 400 \\
\hline
\textbf{รวม} & \textbf{7,811} & \textbf{66} & \textbf{289} & \textbf{6,897} & \textbf{7,951} & \textbf{77} & \textbf{149} & \textbf{6,888} \\
\hline
\end{tabular}
\end{table}

\noindent ผลรวม TP+TN ของ baseline คือ 14,708 และของ PathoWhisper คือ 14,839 จาก 15,000 ช่อง ผลรวมของทั้งสี่คอลัมน์มากกว่า 15,000 เท่ากับจำนวนช่องที่สกัดได้ค่าผิด (63 และ 65 ช่อง)

\section*{ข้อผิดพลาดตามหมวดหมู่}

\begin{table}[H]
\centering
\caption{จำนวนช่องที่สกัดผิดของ PathoWhisper จำแนกตามฟิลด์และหมวดหมู่ (หมวดละ 100 กรณี)}
\label{tab:caterr}
\scriptsize\setlength{\tabcolsep}{2.5pt}
\begin{tabular}{|l|r|r|r|r|r|r|r|r|r|r|r|}
\hline
\textbf{ฟิลด์} & \textbf{1} & \textbf{2} & \textbf{3} & \textbf{4} & \textbf{5} & \textbf{6} & \textbf{7} & \textbf{8} & \textbf{9} & \textbf{10} & \textbf{รวม} \\
\hline
\texttt{s0\_surgical\_no} & 0 & 9 & 0 & 2 & 0 & 0 & 7 & 2 & 0 & 2 & 22 \\
\hline
\texttt{s1\_side} & 0 & 0 & 1 & 0 & 0 & 0 & 0 & 0 & 0 & 0 & 1 \\
\hline
\texttt{s2\_proc} & 0 & 0 & 0 & 0 & 0 & 0 & 0 & 0 & 0 & 0 & 0 \\
\hline
\texttt{s3\_dims} & 1 & 2 & 1 & 4 & 1 & 2 & 5 & 2 & 3 & 2 & 23 \\
\hline
\texttt{s4\_skin} & 0 & 0 & 0 & 0 & 0 & 0 & 2 & 0 & 0 & 0 & 2 \\
\hline
\texttt{s5\_dims} & 0 & 0 & 0 & 0 & 0 & 0 & 7 & 0 & 0 & 0 & 7 \\
\hline
\texttt{s6\_nipple} & 0 & 0 & 0 & 0 & 0 & 0 & 0 & 0 & 0 & 0 & 0 \\
\hline
\texttt{s7\_biopsy\_scar} & 0 & 0 & 0 & 0 & 0 & 0 & 0 & 0 & 0 & 0 & 0 \\
\hline
\texttt{s8\_cavity} & 0 & 0 & 0 & 0 & 0 & 0 & 0 & 0 & 0 & 0 & 0 \\
\hline
\texttt{s9\_residual\_mass} & 0 & 0 & 0 & 0 & 0 & 0 & 0 & 0 & 0 & 0 & 0 \\
\hline
\texttt{s10\_infiltrative} & 0 & 0 & 0 & 1 & 0 & 12 & 5 & 0 & 0 & 0 & 18 \\
\hline
\texttt{s10\_inf\_dims} & 2 & 10 & 5 & 7 & 0 & 0 & 13 & 2 & 1 & 0 & 40 \\
\hline
\texttt{s10\_5\_quadrant} & 0 & 0 & 0 & 2 & 0 & 0 & 0 & 0 & 0 & 0 & 2 \\
\hline
\texttt{s11\_deep\_margin} & 25 & 6 & 0 & 4 & 0 & 0 & 9 & 0 & 0 & 0 & 44 \\
\hline
\texttt{s14\_check} & 0 & 0 & 0 & 0 & 0 & 0 & 2 & 0 & 0 & 0 & 2 \\
\hline
\textbf{รวม} & \textbf{28} & \textbf{27} & \textbf{7} & \textbf{20} & \textbf{1} & \textbf{14} & \textbf{50} & \textbf{6} & \textbf{4} & \textbf{4} & \textbf{161} \\
\hline
\end{tabular}
\end{table}

\noindent หมวดที่ 1--10 เรียงตามตารางที่ \ref{tab:categories} ผลรวม 161 ช่องตรงกับจำนวนช่องผิดของ PathoWhisper ในตารางที่ \ref{tab:confmat}

\section*{การทดสอบหน่วยของตัวปรับข้อความ}

\begin{table}[H]
\centering
\caption{ผลการทดสอบหน่วยของตัวปรับข้อความ (16 กรณี) ผลลัพธ์ที่ได้ตรงกับผลลัพธ์ที่คาดหวังทุกกรณี}
\label{tab:unit}
\scriptsize\setlength{\tabcolsep}{3pt}
\begin{tabular}{|p{1.0cm}|p{5.3cm}|p{5.3cm}|p{1.2cm}|}
\hline
\textbf{ที่} & \textbf{อินพุต} & \textbf{ผลลัพธ์} & \textbf{สถานะ} \\
\hline
TC-01 & infiltrative mass measuring 2.0 x 3.0 x 1.0 cm sorry 2.5 x 3.5 x 1.5 cm & infiltrative mass measuring 2.5 x 3.5 x 1.5 cm & ผ่าน \\
\hline
TC-02 & 10 x 5 x 2 cm, weight is 450 grams & 10 x 5 x 2 cm, weight is 450 grams & ผ่าน \\
\hline
TC-03 & 10 x 5 x 2 cm, weight of 450 grams & 10 x 5 x 2 cm, weight of 450 grams & ผ่าน \\
\hline
TC-04 & 10 x 5 x 2 cm, weight: 450 g & 10 x 5 x 2 cm, weight: 450 g & ผ่าน \\
\hline
TC-05 & 10 x 5 x 2 cm, weight 0.45 kg & 10 x 5 x 2 cm, weight 0.45 kg & ผ่าน \\
\hline
TC-06 & specimen measuring 10 x 5 x 2 cm, weight 450 grams & specimen measuring 10 x 5 x 2 cm, weight 450 grams & ผ่าน \\
\hline
TC-07 & 10 x 5 x 2 cm. actually the skin ellipse measures 5 x 2 cm & 10 x 5 x 2 cm. actually the skin ellipse measures 5 x 2 cm & ผ่าน \\
\hline
TC-08 & 5.3 x 2.4 x 2.9 cm weight margin is 3.0 cm & 5.3 x 2.4 x 2.9 cm weight margin is 3.0 cm & ผ่าน \\
\hline
TC-09 & margin 2 mm no 3 lymph nodes & margin 2 mm no 3 lymph nodes & ผ่าน \\
\hline
TC-10 & mass measuring 2 cm no 3 cm & mass measuring 3 cm & ผ่าน \\
\hline
TC-11 & mass measuring 2.0 x 3.0 x 1.0 cm weight 2.5 x 3.5 x 1.5 cm & mass measuring 2.5 x 3.5 x 1.5 cm & ผ่าน \\
\hline
TC-12 & specimen weight 450 grams & specimen weight 450 grams & ผ่าน \\
\hline
TC-13 & no residual mass identified in the specimen & no residual mass identified in the specimen & ผ่าน \\
\hline
TC-14 & no discrete mass identified & no discrete mass identified & ผ่าน \\
\hline
TC-15 & actually there is no mass in the upper outer quadrant & actually there is no mass in the upper outer quadrant & ผ่าน \\
\hline
TC-16 & mass measuring 2 x 2 cm แก้เป็น 3 x 3 cm & mass measuring 3 x 3 cm & ผ่าน \\
\hline
\end{tabular}
\end{table}

\appendiceschapter{ข้อมูลและสคริปต์ที่ใช้คำนวณผล}

ซอร์สโค้ดของระบบ ตัวสกัดข้อมูล สคริปต์ประเมินผล และชุดข้อความทดสอบทั้งหมดเผยแพร่ในคลังโค้ดสาธารณะ \cite{pathowhisperrepo} โดยไฟล์ที่ใช้คำนวณตัวเลขในเล่มมีรายการดังตารางที่ \ref{tab:appfiles} และคำสั่งทำซ้ำอยู่ท้ายภาคผนวกนี้
\section*{ไฟล์ผลการทดลอง}

\begin{table}[H]
\centering
\caption{ไฟล์ผลการทดลองและสคริปต์}
\label{tab:appfiles}
\scriptsize\setlength{\tabcolsep}{3pt}
\begin{tabular}{|p{6.4cm}|p{6.6cm}|}
\hline
\textbf{ไฟล์} & \textbf{เนื้อหา} \\
\hline
\texttt{benchmark\_1000\_cases\_overnight.csv} & ผลรายกรณี 2,000 บรรทัด (เฉลย ผลถอด WER CER เวลา) \\
\hline
\texttt{benchmark\_1000\_category\_breakdown.csv} & สรุป WER CER เวลา รายหมวด 10 หมวด \\
\hline
\texttt{field\_metrics\_15\_pure\_gt\_v2.csv} & Precision, Recall, F1 และ TP, FP, FN, TN ต่อฟิลด์ (15 ฟิลด์) ของทั้งสองระบบ \\
\hline
\texttt{category\_errors\_15\_fields\_v2.csv} & จำนวนช่องผิดต่อฟิลด์ต่อหมวด \\
\hline
\texttt{parser\_refinement\_comparison.csv} & ผลก่อนและหลังปรับตัวสกัดสองรอบ \\
\hline
\texttt{analyze\_output.txt} & ผลของ \texttt{analyze\_benchmark.py}: สถิติแบบจับคู่ (Wilcoxon, Bootstrap) ผลรายหมวด และอัตราความผิดพลาดตามชนิดของคำ \\
\hline
\texttt{core\_output\_v2.txt} & ผลของ \texttt{eval\_clinical\_core\_metrics.py} กับตัวสกัดรุ่นที่ 2: ตัวชี้วัดระดับเคส และการเปรียบเทียบ 20 กรณี \\
\hline
\texttt{fields\_output\_v2.txt} & ผลของ \texttt{eval\_fields\_15.py}: Macro-F1 ความสอดคล้องระดับช่อง และข้อผิดพลาดต่อฟิลด์ \\
\hline
\texttt{fresh\_testset.json}, \texttt{fresh\_text\_eval.txt} & ชุดข้อความทดสอบใหม่ 300 กรณีและผลการประเมินระดับข้อความ \\
\hline
\texttt{fair\_empirical\_comparison\_20cases.csv} & ผลรายกรณีของการเปรียบเทียบ Vosk, wav2vec 2.0, Whisper Small และ PathoWhisper บน 20 กรณี \\
\hline
\texttt{ablation\_cat8\_fume\_hood.csv} & การเพิ่มองค์ประกอบทีละขั้น (5 กรณี) \\
\hline
\texttt{unit\_test\_results.csv} & ผลการทดสอบหน่วย 16 กรณี \\
\hline
\texttt{web\_functional\_test\_results.csv} & ผลการทดสอบระบบเว็บ 9 โมดูล \\
\hline
\texttt{gt\_audit\_results.csv} & ตรวจว่าเฉลยตรงกับสคริปต์สร้างข้อความ (ไม่ใช่การตรวจความถูกต้องทางคลินิก) \\
\hline
\texttt{ground\_truth\_1000\_pure.json} & เฉลย 15 ฟิลด์ \\
\hline
\texttt{progressive\_ablation\_results.csv}, \texttt{ablation\_summary\_table.csv} & ผลรายเคสและสรุปของการทดลอง Ablation \\
\hline
\texttt{clinical\_case\_level\_evidence.txt} & ตัวอย่างระดับเคสของข้อผิดพลาดหมวด 2 (ตัวสกัดรุ่นแรก) หมวด 6 และด้านของเต้านม \\
\hline
\texttt{analyze\_benchmark.py} & คำนวณสถิติ WER, CER, เวลา และการวิเคราะห์ตามชนิดของคำ \\
\hline
\texttt{eval\_clinical\_core\_metrics.py} & คำนวณ $\mathrm{CER}_{\mathrm{case}}$, Omission/Commission, Flag Coverage, ConER และการทดสอบแบบจับคู่ (ฉบับแก้ไขให้ใช้ตัวเทียบฟิลด์เดียวกับ \texttt{eval\_pure\_gt\_15\_fields.py}) \\
\hline
\texttt{extractor.py}, \texttt{normalizer.py} & ตัวสกัด (รุ่นที่ 2 ที่จัดป้ายมิติตามบริบท) และตัวปรับข้อความ (ใช้ทำซ้ำตัวชี้วัด 15 ฟิลด์) \\
\hline
\texttt{eval\_fields\_15.py} & คำนวณ TP/FP/FN/TN, F1 ต่อฟิลด์ ความสอดคล้องระดับช่อง และข้อผิดพลาดตามหมวด โดยใช้ตัวเทียบเดียวกับ \texttt{eval\_clinical\_core\_metrics.py} \\
\hline
\texttt{test\_extractor\_context.py} & ทดสอบการกำหนดมิติตามบริบทและกรณี skin ellipse 7 กรณี \\
\hline
\texttt{generate\_fresh\_testset.py}, \texttt{eval\_fresh.py} & สร้างและประเมินชุดข้อความทดสอบใหม่ 300 กรณี \\
\hline
\texttt{eval\_pure\_gt\_15\_fields.py} & สร้างเฉลยและคำนวณตัวชี้วัดการสกัด 15 ฟิลด์ \\
\hline
\texttt{gt\_audit.py} & ตรวจว่าเฉลยตรงกับสคริปต์สร้างข้อความ \\
\hline
\texttt{run\_unit\_tests.py}, \texttt{test\_web\_full.py} & ทดสอบหน่วยและทดสอบระบบเว็บ \\
\hline
\texttt{diagnose\_case\_evidence.py}, \texttt{generate\_publication\_charts.py} & สคริปต์วิเคราะห์ข้อผิดพลาดระดับเคสและสร้างแผนภาพ \\
\hline
\end{tabular}
\end{table}

\noindent หมายเหตุการทำซ้ำ: ตัวเลขทั้งหมดในเล่มรันในสภาพแวดล้อมที่ติดตั้ง spaCy (\texttt{en\_core\_web\_sm} 3.8.0) เมื่อรันซ้ำโดยไม่ติดตั้ง spaCy ค่า TP, FN และ TN ตรงกับตารางทุกฟิลด์ (ความสอดคล้องระดับช่อง 14,708 และ 14,839) และ $\mathrm{CER}_{\mathrm{case}}$ เท่ากัน (ร้อยละ 18.80 และ 11.10) ต่างที่จำนวน FP ของ \texttt{s3\_dims} (3 และ 4 ช่อง แทน 12 และ 14) และของ \texttt{s11\_deep\_margin} ของ baseline (0 แทน 1 ช่อง) การเว้นว่าง/การกรอกค่าผิดรวม (236/56 และ 94/67 แทน 226/66 และ 84/77) และสัดส่วนเคสวิกฤตผิดที่มีธงเตือน (ร้อยละ 18.1 และ 36.9 แทน 13.8 และ 30.6) ซึ่งสอดคล้องกับการที่กลไกสำรองของ spaCy เปลี่ยนการเว้นว่างบางส่วนเป็นค่าผิด

\begin{verbatim}
pip install pandas scipy jiwer
python eval_pure_gt_15_fields.py
python scripts/eval_clinical_core_metrics.py \
    <gt.json> <benchmark.csv> <fair_20cases.csv>
python scripts/eval_fields_15.py \
    <gt.json> <benchmark.csv> <outdir>
python test_extractor_context.py
\end{verbatim}

\noindent สคริปต์ \texttt{analyze\_benchmark.py} ใช้คำนวณสถิติ WER ทั้งเกณฑ์ A และเกณฑ์ B, CER, เวลาประมวลผล, ช่วงความเชื่อมั่นแบบ Bootstrap (5,000 รอบ ที่ระดับความเชื่อมั่น 95\%), การทดสอบ Wilcoxon Signed-Rank ตลอดจนการวิเคราะห์ความผิดพลาดจำแนกตามหมวดหมู่และชนิดของคำศัพท์จากข้อมูลผลการทดสอบรายกรณี

\appendiceschapter{หลักฐานการทดสอบการทำงานแบบออฟไลน์}

การทดสอบระบบในสภาวะตัดการเชื่อมต่อเครือข่ายภายนอก (Isolated Environment Testing) ดำเนินการบนเครื่องคอมพิวเตอร์ที่ใช้ในการทดลอง เพื่อยืนยันว่าระบบ PathoWhisper สามารถประมวลผลสัญญาณเสียง สกัดข้อมูลลงแบบฟอร์ม และสร้างรายงานสรุปในรูปแบบเอกสาร PDF ได้อย่างสมบูรณ์โดยไม่พึ่งพาการเชื่อมต่ออินเทอร์เน็ตหรือบริการภายนอก ซึ่งสอดคล้องกับมาตรการคุ้มครองข้อมูลส่วนบุคคลทางการแพทย์ ขั้นตอนและผลการทดสอบเชิงฟังก์ชันสรุปดังตารางที่ \ref{tab:offline}

\begin{table}[H]
\centering
\caption{สรุปขั้นตอนและผลการทดสอบการทำงานของระบบในสภาวะออฟไลน์}
\label{tab:offline}
\footnotesize\setlength{\tabcolsep}{4pt}
\begin{tabular}{|c|p{4.2cm}|p{7.8cm}|}
\hline
\textbf{ลำดับ} & \textbf{ขั้นตอนการทดสอบ} & \textbf{ผลการทดสอบและหลักฐานเชิงประจักษ์} \\
\hline
1 & การตรวจสอบสถานะเครือข่าย (ภาพที่ \ref{fig:off1}) & คำสั่ง \texttt{ping 8.8.8.8} ล้มเหลวทั้ง 4 ครั้ง (ส่ง 4 รับ 0 สูญหายร้อยละ 100) และไอคอนเครือข่ายบนแถบงานแสดงว่าไม่มีอินเทอร์เน็ต \\
\hline
2 & การประมวลผลเสียงและสกัดข้อมูล (ภาพที่ \ref{fig:off2}) & ระบบถอดความสัญญาณเสียงความยาว 23 วินาที และสกัดข้อมูลลงแบบฟอร์มได้ครบถ้วน เช่น เลขที่สิ่งส่งตรวจ S-24-1027, ประเภทหัตถการ MRM ข้างขวา, ขนาดสิ่งส่งตรวจ และตำแหน่งก้อน \\
\hline
3 & การตรวจสอบเงื่อนไขความถูกต้อง (ภาพที่ \ref{fig:off3}) & เมื่อทดสอบด้วยกรณีที่ไม่มีการระบุเลขที่สิ่งส่งตรวจ ระบบแสดงข้อความแจ้งเตือนและเว้นว่างฟิลด์ดังกล่าวตามเงื่อนไขความปลอดภัยทางข้อมูล \\
\hline
4 & การสร้างรายงานสรุปผล PDF (ภาพที่ \ref{fig:off4}) & ไฟล์รายงาน \texttt{pathology\_breast\_gross\_S-24-1027 (1).pdf} ของเคสเดียวกับภาพที่ \ref{fig:off2} เปิดขณะที่ไอคอนเครือข่ายยังแสดงว่าไม่มีอินเทอร์เน็ต ค่าที่ปรากฏตรงกับที่ผู้พูดกล่าว \\
\hline
\end{tabular}
\end{table}

\begin{figure}[H]
\centering
\includegraphics[width=0.97\textwidth]{images/evidence/offline_1_ping_1553.png}
\caption{การตรวจสอบสถานะการตัดการเชื่อมต่อเครือข่ายภายนอก (Network Isolation Verification)}
\label{fig:off1}
\end{figure}

\begin{figure}[H]
\centering
\includegraphics[width=0.97\textwidth]{images/evidence/offline_2_dictation_1742.png}
\caption{การประมวลผลถอดความเสียงและสกัดข้อมูลลงแบบฟอร์มในสภาวะออฟไลน์}
\label{fig:off2}
\end{figure}

\begin{figure}[H]
\centering
\includegraphics[width=0.97\textwidth]{images/evidence/offline_3_missing_field_flag_1600.png}
\caption{การแสดงข้อความแจ้งเตือนของระบบเมื่อไม่พบข้อมูลเลขที่สิ่งส่งตรวจทางศัลยกรรม}
\label{fig:off3}
\end{figure}

\begin{figure}[H]
\centering
\includegraphics[width=0.97\textwidth]{images/evidence/offline_4_pdf_1742.png}
\caption{รายงาน PDF ของเคส S-24-1027 ที่เปิดขณะเครือข่ายถูกตัดการเชื่อมต่อ}
\label{fig:off4}
\end{figure}

\noindent \textbf{ข้อสังเกตเชิงเทคนิคจากการทดสอบ:}
\begin{itemize}
    \item \textbf{ฟิลด์หัวนม:} ผู้พูดกล่าวว่า \textit{nipple is retracted} แต่ช่องหัวนมในแบบฟอร์มไม่ถูกเลือก (ภาพที่ \ref{fig:off2}) สอดคล้องกับข้อจำกัดว่าฟิลด์หัวนมไม่ได้ผ่านการทดสอบการตรวจจับในชุดทดสอบ
    \item \textbf{ขนาดผิวหนังกับระยะห่างขอบตัด:} ในรายงาน PDF ช่อง \textit{cm. from skin} ว่าง (ภาพที่ \ref{fig:off4}) ทั้งที่ผู้พูดกล่าวขนาดวงรีผิวหนัง 7 x 3 cm การทดสอบรอบแรกพบว่าตัวสกัดเคยกรอกค่า 7 ในช่องนี้ จึงแก้กฎให้ต้องมีคำว่า margin และสร้างโปรแกรมใหม่ก่อนถ่ายภาพนี้
    \item \textbf{การจัดสรรชุดมิติตามบริบท:} ผู้พูดบอกขนาดก้อนก่อนขนาดสิ่งส่งตรวจ แต่ตัวสกัดกรอกขนาดทั้งสองถูกช่อง (ภาพที่ \ref{fig:off2}) การทดสอบนี้เป็นการสาธิตบนเคสเดียว มิใช่การวัดความแม่นยำ
\end{itemize}